% =====================================================================
% cabezazos.tex — Log de cabezazos (todas las fases)
% Requiere estilo.tex. Las filas de la semana 2 enlazan con los cuadros
% detallados de semana2.tex.
% =====================================================================

\chapter{Log de cabezazos}
\label{cap:cabezazos}

Registro de los problemas encontrados durante el despliegue, su causa y cómo se
resolvieron. Los de la semana 2 se explican con más detalle, junto con sus
capturas, en el capítulo~\ref{cap:semana2}.

{\footnotesize
\let\cmd\cmdb
\renewcommand{\arraystretch}{1.25}
\begin{xltabular}{\linewidth}{@{}>{\raggedright\arraybackslash}p{0.4cm}
                              >{\raggedright\arraybackslash}p{1.5cm}
                              >{\raggedright\arraybackslash}p{1.5cm}
                              >{\raggedright\arraybackslash}X
                              >{\raggedright\arraybackslash}X
                              >{\raggedright\arraybackslash}X@{}}
\toprule
\textbf{Nº} & \textbf{Fecha} & \textbf{Fase} & \textbf{Problema} & \textbf{Causa} & \textbf{Solución} \\
\midrule
\endfirsthead
\toprule
\textbf{Nº} & \textbf{Fecha} & \textbf{Fase} & \textbf{Problema} & \textbf{Causa} & \textbf{Solución} \\
\midrule
\endhead
\midrule
\multicolumn{6}{r}{\textit{Continúa en la página siguiente}} \\
\endfoot
\bottomrule
\endlastfoot

\rowcolor{gray!12}\multicolumn{6}{@{}l}{\textbf{Semana 1: máquina virtual, red y preparación OFA}} \\

1 & 06/10/2026 & VM / red &
\cmd{virt-install} falla: «No se puede obtener MTU del interfaz en 'br0': No existe el dispositivo». &
El bridge \cmd{br0} no existía en el anfitrión. Después de crearlo con
\cmd{nmcli con add}, el error se repitió porque la conexión estaba definida pero
no activada. &
Activar el bridge con
\texttt{nmcli con down "Conexión cableada 1"; nmcli con up br0-slave; nmcli con up br0},
comprobar con \cmd{ip -br link} y desactivar el \emph{autoconnect} de la
conexión antigua. Solo funciona con cable, no con Wi-Fi. \\

2 & 06/10/2026 & VM &
\cmd{virsh shutdown delfos} devuelve «error: incorrecto al obtener dominio 'delfos'». &
La VM se creó en la conexión del sistema (\cmd{qemu:///system}), pero
\cmd{virsh} sin \cmd{sudo} usa por defecto la de sesión del usuario
(\cmd{qemu:///session}), que no la ve. &
Usar \cmd{virsh -c qemu:///system ...} o fijar
\cmd{export LIBVIRT_DEFAULT_URI="qemu:///system"} en el \cmd{.bashrc}. \\

3 & 06/10/2026 & OFA &
No existe en Ubuntu el paquete \cmd{oracle-database-preinstall} que prepara el
sistema; \cmd{/etc/sysctl.d} solo contiene un README. &
Oracle solo publica ese paquete para Oracle Linux / RHEL. &
Crear a mano \cmd{/etc/sysctl.d/99-oracle.conf} (parámetros del kernel),
\cmd{/etc/security/limits.d/oracle.conf} (límites), el usuario \cmd{oracle},
los grupos y los directorios OFA. \\

\rowcolor{gray!12}\multicolumn{6}{@{}l}{\textbf{Semana 2: Oracle XE 21c y BAHIA.inf}} \\

\ref{cab:disco} & 13/10/2026 & Instalación XE &
\cmd{alien} aborta con «No space left on device». &
\cmd{alien} necesita $\sim$20 GB temporales y el volumen lógico raíz solo tenía
19 GB (el instalador dejó otros 19 GB sin asignar en el VG). &
\texttt{lvextend -r -l +100\%FREE}; además \cmd{qemu-img resize +30G},
\cmd{growpart}, \cmd{pvresize} y \cmd{lvextend}: raíz de 67 GB. \\

\ref{cab:tmpfs} & 13/10/2026 & Instalación XE &
«zstd write error: No space left on device» con 55 GB libres; el
\cmd{.deb} sale de 521 KB. &
\cmd{dpkg-deb} comprime en \cmd{/tmp}, que en Ubuntu 26.04 es un tmpfs en RAM de
1,7 GB. &
\cmd{sudo env TMPDIR=/home/sercal/tmp-alien alien --scripts -d ...}. Los errores
de \cmd{dpkg-shlibdeps} se ignoran. \\

\ref{cab:awk} & 13/10/2026 & Instalación XE &
El \emph{preinst} falla con «awk: command not found». &
\cmd{ln -sf /usr/bin/awk /bin/awk} en un sistema con \emph{usrmerge} creó el
enlace circular \cmd{/usr/bin/awk -> /usr/bin/awk}. &
\cmd{sudo rm /usr/bin/awk} y \cmd{sudo apt install --reinstall mawk}. No crear
enlaces para \cmd{awk}. \\

\ref{cab:oraclebase} & 13/10/2026 & Instalación XE &
«ORACLE\_BASE directory (/opt/oracle) is not owned by oracle user». &
El \emph{preinst} comprueba la propiedad de \cmd{/opt/oracle} antes de
desempaquetar. &
\cmd{chown -R oracle:oinstall /opt/oracle /opt/oraInventory} antes de instalar
(nunca después). \\

\ref{cab:display} & 13/10/2026 & DBCA &
\cmd{dbca}: «DISPLAY not set». &
Sesión abierta con \cmd{sudo su - oracle} sin reenvío X11; al pegar un bloque,
las líneas tras \cmd{sudo su} se pierden. &
Conectar con \cmd{ssh -Y oracle@delfos.inf} o copiar la cookie con
\cmd{xauth add}, ejecutando las órdenes de una en una. \\

\ref{cab:cookie} & 13/10/2026 & DBCA &
«Invalid MIT-MAGIC-COOKIE-1 key» y
\cmd{NoClassDefFoundError: ... XToolkit}. &
El escritorio del anfitrión (Wayland) no da una cookie X11 válida:
«No xauth data; using fake authentication data». &
Usar DBCA en modo silencioso (\cmd{-silent}). \\

\ref{cab:dbt12513} & 13/10/2026 & DBCA &
\cmd{[DBT-12513]} El esquema de muestra no está soportado. &
DBCA 21c ya no incluye los esquemas de ejemplo. &
Crear sin \cmd{-sampleSchema} e instalarlos desde
\cmd{oracle-samples/db-sample-schemas}. \\

\ref{cab:npe} & 13/10/2026 & DBCA &
\cmd{-generateScripts} acaba con \cmd{java.lang.NullPointerException}. &
Fallo de DBCA al terminar; el progreso es simulado y no se crea nada. &
Ninguna: los 17 scripts se generan completos. Verificar que no hay
\cmd{pmon} ni ficheros en \cmd{oradata}. \\

\ref{cab:oradism} & 13/10/2026 & DBCA &
\cmd{[DBT-10011]} propiedad/permisos incorrectos de \cmd{oradism}
(\cmd{PRVG-11960}, \cmd{PRVG-2031}). &
Un \cmd{chown -R} posterior a la instalación quitó a \cmd{root} la propiedad y
el bit SUID de \cmd{oradism}. &
\cmd{sudo $ORACLE_HOME/root.sh}. No volver a usar \cmd{chown -R} sobre
\cmd{/opt/oracle}. \\

\ref{cab:pwd} & 13/10/2026 & DBCA &
Aviso \cmd{[DBT-06208]} para SYS, SYSTEM y PDBADMIN. &
Contraseñas cortas o prohibidas (p.\,ej. «oracle»). &
Contraseñas de 8+ caracteres con mayúsculas, minúsculas y dígitos, guardadas
fuera del repositorio. \\

\ref{cab:pdb} & 13/10/2026 & BAHIA &
Tras reiniciar, \cmd{BAHIAPDB} queda en \cmd{MOUNTED}. &
Las PDB no se abren solas al arrancar la CDB. &
\cmd{ALTER PLUGGABLE DATABASE BAHIAPDB OPEN;} y \cmd{... SAVE STATE;}. \\

\ref{cab:ora65096} & 13/10/2026 & Esquemas &
\cmd{hr_install.sql}: \cmd{ORA-65096} nombre de usuario o rol común no válido. &
Se ejecutó en \cmd{CDB$ROOT}, donde solo se admiten usuarios comunes
\texttt{C\#\#}. &
Conectar a \cmd{//delfos.inf:1521/bahiapdb.inf} y comprobar
\cmd{SHOW CON_NAME}. \\

\ref{cab:ora01017} & 13/10/2026 & Esquemas &
\cmd{ORA-01017} usuario/contraseña no válidos al conectar como SYSTEM. &
Contraseña mal tecleada. &
Reintroducir las credenciales. \\

\ref{cab:sudo} & 13/10/2026 & Arranque &
«I'm sorry oracle. I'm afraid I can't do that» al usar \cmd{sudo}. &
El usuario \cmd{oracle} no tiene (ni debe tener) privilegios de \cmd{sudo}. &
Ejecutar las tareas de sistema con \cmd{sercal} (\cmd{su - sercal}). \\
\end{xltabular}
}

\section*{Lecciones aprendidas}

\begin{itemize}[leftmargin=*]
  \item Comprobar el espacio \emph{real} disponible (\cmd{df -h /}, \cmd{vgs},
        \cmd{df -h /tmp}) antes de procesar paquetes grandes.
  \item No aplicar a ciegas recetas escritas para distribuciones antiguas: en
        Ubuntu moderno \cmd{/bin} y \cmd{/usr/bin} son el mismo directorio.
  \item No usar \cmd{chown -R} sobre un \texttt{ORACLE\_HOME} ya instalado: hay
        binarios que deben ser de \cmd{root} con SUID.
  \item En una CDB, los esquemas de aplicación van siempre en una PDB.
  \item Separar los usuarios: \cmd{sercal} administra el sistema y \cmd{oracle}
        administra la base de datos.
  \item Hacer instantáneas antes de cada paso delicado ha permitido trabajar sin
        miedo a tener que reinstalar desde cero.
\end{itemize}
